\documentclass[11pt,a4paper]{article}

\usepackage[T1]{fontenc}
\usepackage[utf8]{inputenc}
\usepackage[brazil]{babel}
\usepackage{amsmath,amssymb}
\usepackage{booktabs}
\usepackage{geometry}
\usepackage{hyperref}
\usepackage{listings}
\usepackage{longtable}
\usepackage{microtype}
\usepackage{tabularx}
\usepackage{xcolor}

\geometry{top=2.2cm,bottom=2.2cm,left=2.2cm,right=2.2cm}
\hypersetup{colorlinks=true,linkcolor=blue!45!black,urlcolor=blue!45!black}
\definecolor{navy}{HTML}{032147}
\definecolor{blue}{HTML}{209DD7}
\definecolor{purple}{HTML}{753991}
\definecolor{codebg}{HTML}{F3F6F9}
\definecolor{codeborder}{HTML}{D8E1EA}
\definecolor{muted}{HTML}{52606D}
\lstset{
  language=Python,
  basicstyle=\ttfamily\small,
  backgroundcolor=\color{codebg},
  frame=single,
  rulecolor=\color{codeborder},
  numbers=left,
  numberstyle=\tiny\color{muted},
  breaklines=true,
  showstringspaces=false,
  columns=fullflexible,
  keywordstyle=\color{purple}\bfseries,
  commentstyle=\color{muted}
}

\title{\textbf{Análise de Complexidade Big-O}\\[0.4em]
\large Estruturas centrais e A* dependente do tempo}
\author{Ferramenta de Otimização de Rotas de Transporte Público}
\date{20 de novembro de 2026}

\begin{document}
\maketitle

\begin{abstract}
Este relatório fundamenta as complexidades temporal e espacial das estruturas
centrais do projeto a partir do código efetivamente executado. A representação
por listas de adjacência ocupa $O(V+E)$, adiciona vértices e
recupera adjacências em $O(1)$ médio, mas a inserção ordenada de uma conexão
custa $O(K_v)$ no pior caso por causa do deslocamento em listas Python. No
roteamento, a busca binária encontra a primeira partida viável em
$O(\log K_v)$, porém o sufixo relevante ainda é percorrido. Para a implementação
com \texttt{heapq} e entradas obsoletas, o limite rigoroso do A* é
$O(E\log(E+1)+V\log(K_{\max}+1)+L+V+1)$ no pior caso médio das tabelas hash e
$O(V+E)$ de espaço auxiliar. Também são declaradas as hipóteses que sustentam a
admissibilidade e a consistência da heurística de Haversine.
\end{abstract}

\section{Escopo, modelo e notação}

O sistema modela uma rede dirigida e temporal. Uma estação é um vértice e cada
\texttt{Connection} é uma partida específica, com destino, horário e duração.
Assim, duas viagens entre o mesmo par de estações em horários diferentes são
arestas distintas. A análise usa:

\begin{itemize}
  \item $V$: número de estações registradas;
  \item $E$: número total de conexões temporais;
  \item $K_v$: tamanho da lista de conexões de saída da estação $v$;
  \item $K_{\max}=\max_v K_v$, definido como zero se não houver estações;
  \item $X$: conjunto de estações efetivamente expandidas pelo A*;
  \item $H$: número de inserções no heap durante uma consulta;
  \item $L$: número de trechos da rota reconstruída.
\end{itemize}

As operações em \texttt{dict} e \texttt{set} são tratadas como $O(1)$ em
tempo médio amortizado, conforme o modelo usual das tabelas hash. Colisões
patológicas podem degradá-las a $O(V)$, mas não representam o comportamento
esperado com identificadores de estação imutáveis. O custo de calcular o hash
do identificador é considerado constante em relação a $V$ e $E$.

\section{Representação do grafo e consumo de espaço}

\texttt{TransportGraph} mantém um único dicionário de adjacência. Cada chave
referencia uma estação e cada valor referencia uma lista ordenada de conexões.

\begin{lstlisting}[caption={Estrutura e operações básicas em \texttt{backend/app/graph.py:9--32}.}]
class TransportGraph:
    __slots__ = ("_adjacency_list",)

    def __init__(self) -> None:
        self._adjacency_list: dict[Station, list[Connection]] = {}

    def add_station(self, station: Station) -> None:
        self._adjacency_list.setdefault(station, [])

    def get_neighbors(self, station: Station) -> list[Connection]:
        return self._adjacency_list[station]
\end{lstlisting}

Há uma entrada por vértice e uma referência em alguma lista para cada aresta.
Logo, a estrutura ocupa $O(V+E)$. As classes de domínio usam
\texttt{slots=True}, eliminando o dicionário de atributos por instância. Isso
reduz constantes de memória, mas não altera a ordem assintótica: cada
\texttt{Station} e cada \texttt{Connection} continuam ocupando $O(1)$.

\section{Inserções e consultas no \texttt{TransportGraph}}

\subsection{Adição de vértice}

\texttt{setdefault} executa uma busca hash e só cria a lista vazia se a estação
ainda não existir. A complexidade é $O(1)$ média e amortizada; uma realocação
ocasional do dicionário custa $O(V)$, distribuída entre inserções sucessivas.
A memória incremental é $O(1)$ por vértice.

\subsection{Adição ordenada de conexão}

\begin{lstlisting}[caption={Inserção por horário em \texttt{backend/app/graph.py:20--29}.}]
connections = self._adjacency_list[origin]
if connection.destination not in self._adjacency_list:
    raise KeyError(connection.destination)
insertion_index = bisect_right(
    connections,
    connection.departure_time,
    key=lambda item: item.departure_time,
)
connections.insert(insertion_index, connection)
\end{lstlisting}

As duas consultas ao dicionário são $O(1)$ em média. A busca binária consome
$O(\log(K_v+1)+1)$ de tempo, incluindo a chamada sobre lista vazia.
Entretanto, \texttt{list.insert} precisa deslocar
até $K_v$ referências quando a posição está no início da lista. Portanto:

\[
T_{\text{add\_connection}}(v)
=O(1)+O(\log(K_v+1))+O(K_v+1)=O(K_v+1).
\]

Este ponto impede classificar o pior caso da inserção de arestas como
$O(\log(K_v+1)+1)$. Quando a posição é o final, a mutação da lista custa
$O(1)$ amortizado, pois não desloca referências; a operação completa ainda
executa \texttt{bisect\_right} e custa $O(\log(K_v+1)+1)$.
No pior caso, custa $O(K_v+1)$, equivalente a $O(K_v)$ quando $K_v\ge 1$.

\subsection{Recuperação de vizinhos}

\texttt{get\_neighbors} realiza uma única consulta ao dicionário e devolve a
lista interna, sem copiá-la. O tempo é $O(1)$ médio e o espaço adicional é
$O(1)$. Percorrer o resultado é uma operação separada e custa $O(K_v)$.

\begin{table}[ht]
\centering
\caption{Complexidades das operações centrais do grafo.}
\begin{tabularx}{\textwidth}{@{}lXX@{}}
\toprule
Operação & Tempo & Espaço adicional \\
\midrule
\texttt{add\_station} & $O(1)$ médio/amortizado & $O(1)$ \\
\texttt{get\_neighbors} & $O(1)$ médio & $O(1)$ \\
Percorrer vizinhos de $v$ & $O(K_v)$ & $O(1)$ \\
\texttt{add\_connection} & $O(K_v+1)$ pior caso & $O(1)$ \\
\texttt{iter\_stations} & $O(1)$ para criar; $O(V)$ para consumir & $O(1)$ \\
Remover vértice temporário & $O(V+E)$ & $O(K_{\max})$ de pico \\
\bottomrule
\end{tabularx}
\end{table}

\section{Busca temporal de partidas}

As listas são ordenadas por \texttt{departure\_time}; por isso o roteador usa
\texttt{bisect\_left} para encontrar a primeira partida com horário maior ou
igual ao instante de chegada $t$.

\begin{lstlisting}[caption={Busca e varredura em \texttt{backend/app/route\_simulator.py:107--124}.}]
start = bisect_left(connections, t, key=lambda c: c.departure_time)
for i in range(start, len(connections)):
    connection = connections[i]
    if connection.vehicle_type is VehicleType.WALK:
        continue
    if connection.departure_time >= g.get(destination, float("inf")):
        break
    arrival = connection.departure_time + connection.travel_time
    if arrival < g.get(connection.destination, float("inf")):
        # atualiza rotulo, predecessor e heap
\end{lstlisting}

Localizar \texttt{start} custa $O(\log(K_v+1)+1)$. Se $S_v$ conexões do sufixo forem
inspecionadas até o fim ou até a poda, a expansão custa
$O(\log(K_v+1)+S_v+1)$. A busca binária evita examinar partidas já perdidas, mas não
transforma toda a expansão em $O(\log K_v)$. Como cada estação é fechada no
máximo uma vez, o número total de conexões inspecionadas é $O(E)$.

As conexões de caminhada possuem horário zero e formam um prefixo. O primeiro
laço para ao encontrar horário positivo, logo custa $O(W_v+1)$, em que $W_v$ é o
tamanho desse prefixo. Somado sobre as expansões, custa $O(E+|X|)$, incluindo
a verificação do laço mesmo quando a lista está vazia.

\section{A* dependente do tempo}

\subsection{Modelo de custo}

Para uma conexão agendada $c$ embarcada no instante $t$,

\[
c_t(n,m)=(\text{departure}_c-t)+\text{travel}_c.
\]

O rótulo $g[v]$ é o horário absoluto de chegada. Caminhadas partem
imediatamente e usam $c_t(n,m)=\text{travel}_c$. Sob a propriedade FIFO, chegar
mais tarde a uma estação não permite chegar antes pela mesma conexão. Assim, o
algoritmo pode fechar cada estação uma única vez.

\subsection{Heurística, admissibilidade e consistência}

A implementação calcula

\[
h(n)=\frac{d_H(n,d)}{V_{\max}},
\]

onde $d_H$ é Haversine e $V_{\max}$ é a maior razão
\texttt{distance/travel\_time} observada. Haversine executa quantidade fixa de
operações trigonométricas, portanto $h(n)$ custa $O(1)$.

Pela desigualdade triangular,
$d_H(n,d)\le d_H(n,m)+d_H(m,d)$. Como a distância das conexões carregadas é
calculada por Haversine e $V_{\max}$ limita a velocidade,

\[
\frac{d_H(n,m)}{V_{\max}}
\le \text{travel}(n,m)
\le c_t(n,m).
\]

Somando $h(m)$, obtém-se
$h(n)\le c_t(n,m)+h(m)$ para toda aresta válida. Portanto, a heurística é
consistente e, consequentemente, admissível. A garantia depende de
\texttt{max\_speed} realmente limitar todas as conexões e de
\texttt{distance} não subestimar a distância geodésica.

\subsection{Fila de prioridade e relaxamentos}

\begin{lstlisting}[caption={Núcleo da fila em \texttt{backend/app/route\_simulator.py:67--82}.}]
heap = [
    (departure_time + self.heuristic(origin, destination),
     next(counter), origin)
]
while heap:
    _, _, station = heappop(heap)
    if station in closed:
        continue
    if station == destination:
        return self._reconstruct_path(predecessor, origin, destination)
    closed.add(station)
\end{lstlisting}

\texttt{heapq} não oferece \emph{decrease-key}. Quando um rótulo melhora, uma
nova tupla é inserida e a antiga permanece até ser removida como obsoleta. Se
$H$ é o número de inserções, há no máximo $H$ remoções; cada operação custa
$O(\log(H+1))$. Cada relaxamento bem-sucedido corresponde a uma conexão examinada,
então $H\le O(E+1)$.

\subsection{Complexidade temporal rigorosa}

Separando os custos da implementação:

\[
T=O\left(
E+\sum_{v\in X}\log(K_v+1)+H\log(H+1)+L+|X|+1
\right).
\]

Com $|X|\le V$, $H=O(E+1)$ e $K_v\le K_{\max}$:

\[
T=O\left(E\log(E+1)+V\log(K_{\max}+1)+L+V+1\right).
\]

Os termos $|X|+1$ e $V+1$ explicitam a inicialização e o custo constante por
expansão omitidos na forma anterior aprovada. Com $E=K_{\max}=L=0$, uma busca
entre estações isoladas ainda inicializa e remove a entrada da origem no heap.
A nova expressão cobre esse caso sem atribuir trabalho zero. No caso usual em
que o termo do heap domina, escreve-se $O(E\log(E+1))$.
A forma clássica $O((V+E)\log V)$ pressupõe heap limitado a $O(V)$ ou uma
operação de \emph{decrease-key}. Como esta implementação mantém duplicatas
obsoletas, $\log E$ é o limite mais fiel. A heurística pode reduzir $|X|$, $H$ e
as varreduras na prática, mas não melhora o pior caso assintótico: uma
heurística nula ou pouco informativa aproxima o comportamento de Dijkstra.

\subsection{Complexidade espacial rigorosa}

Os mapas \texttt{g} e \texttt{predecessor} e o conjunto \texttt{closed}
ocupam $O(V)$. O heap pode acumular $O(E+1)$ entradas por causa da entrada inicial,
das melhorias de
rótulo e das entradas obsoletas. A rota retornada ocupa $O(L)$. Portanto, além
do grafo já carregado:

\[
S_{A^*}=O(V+E+L+1)=O(V+E+1).
\]

Para uma rede não vazia, o espaço total equivale a $O(V+E)$. Classificar o heap
como $O(V)$ seria correto apenas com uma entrada ativa por
vértice, condição não garantida pelo uso atual de \texttt{heapq}.

\section{Operações de apoio}

\begin{longtable}{@{}p{0.31\textwidth}p{0.28\textwidth}p{0.31\textwidth}@{}}
\caption{Custos adicionais relevantes.}\\
\toprule
Operação & Tempo & Justificativa \\
\midrule
\endfirsthead
\toprule
Operação & Tempo & Justificativa \\
\midrule
\endhead
\texttt{max\_speed\_of} & $O(V+E)$ & Visita cada estação e cada conexão uma vez. \\
\texttt{heuristic} / Haversine & $O(1)$ & Número fixo de operações aritméticas. \\
Reconstrução do caminho & $O(L)$ & Segue um predecessor por trecho e inverte a lista. \\
Localizar estação para snap & $O(V)$ & Calcula Haversine para cada estação real. \\
Limpar vértice temporário & $O(V+E)$ & Filtra todas as listas de adjacência. \\
\bottomrule
\end{longtable}

\section{Exemplos aplicados ao código}

\paragraph{Exemplo 1 - inserção.}
Se uma estação possui $K_v=200$ partidas, \texttt{bisect\_right} precisa de
cerca de $\log_2 200\approx 8$ comparações. Inserir uma nova caminhada no
início, entretanto, pode deslocar as 200 referências. O custo dominante é
linear em $K_v$.

\paragraph{Exemplo 2 - consulta temporal.}
Após 150 das 200 partidas, a busca \texttt{bisect\_left}
descarta esse prefixo em cerca de oito comparações. Se as 50 partidas restantes
forem relevantes antes da poda, a expansão custa $O(\log 200+50)$, não apenas
$O(\log 200)$.

\paragraph{Exemplo 3 - entradas obsoletas no A*.}
Uma estação pode receber rótulos progressivamente melhores por predecessores
diferentes. Cada melhora adiciona uma tupla ao heap; as anteriores não são
removidas imediatamente. O teste \texttt{if station in closed} descarta essas
tuplas depois, preservando a correção, mas permitindo $O(E)$ itens no heap.

\section{Síntese}

\begin{table}[ht]
\centering
\caption{Limites assintóticos consolidados.}
\begin{tabularx}{\textwidth}{@{}lXX@{}}
\toprule
Componente & Tempo & Espaço \\
\midrule
Grafo completo & -- & $O(V+E)$ \\
Adicionar estação & $O(1)$ médio/amortizado & $O(1)$ incremental \\
Adicionar conexão & $O(K_v+1)$ pior caso & $O(1)$ incremental \\
Obter lista de vizinhos & $O(1)$ médio & $O(1)$ \\
Encontrar primeira partida & $O(\log(K_v+1)+1)$ & $O(1)$ \\
Expandir estação & $O(W_v+\log(K_v+1)+S_v+1)$ & até $O(K_v+1)$ no heap \\
A* - pior caso rigoroso & $O(E\log(E+1)+V\log(K_{\max}+1)+L+V+1)$ & $O(V+E+1)$ \\
Reconstruir rota & $O(L)$ & $O(L)$ \\
\bottomrule
\end{tabularx}
\end{table}

O projeto cumpre a meta de $O(1)$ médio para adicionar vértices e recuperar a
referência dos vizinhos. A ordenação na escrita desloca custo para a carga do
grafo e permite localizar partidas perdidas em tempo logarítmico durante muitas
consultas. O A* permanece ótimo sob FIFO e heurística consistente. Sua análise
deve, porém, contabilizar a varredura após o \texttt{bisect} e as duplicatas do
\texttt{heapq}; ignorá-las produz limites de tempo e espaço excessivamente
otimistas.

\section*{Referências}

\begin{enumerate}
  \item Cormen, T. H.; Leiserson, C. E.; Rivest, R. L.; Stein, C.
  \textit{Introduction to Algorithms}. 4. ed. MIT Press, 2022.
  \item Hart, P. E.; Nilsson, N. J.; Raphael, B. A formal basis for the
  heuristic determination of minimum cost paths. \textit{IEEE Transactions on
  Systems Science and Cybernetics}, v. 4, n. 2, p. 100--107, 1968.
  \item Python Software Foundation. \textit{bisect --- Array bisection
  algorithm}; \textit{heapq --- Heap queue algorithm}; \textit{Data model}.
  Documentação da biblioteca padrão Python, versão 3.
  \item Código-fonte do projeto: \path{backend/app/graph.py},
  \path{backend/app/domain.py}, \path{backend/app/data_loader.py} e
  \path{backend/app/route_simulator.py}, revisão de 20/11/2026.
\end{enumerate}

\end{document}
